\documentclass[11pt]{article}

\usepackage[T1]{fontenc}
\usepackage{lmodern}
\usepackage[margin=1in]{geometry}
\usepackage{amsmath,amssymb,mathtools,bm}
\usepackage{xcolor}
\usepackage{hyperref}

\setlength{\parskip}{0.45em}
\setlength{\parindent}{0pt}
\setlength{\emergencystretch}{2em}
\allowdisplaybreaks

\hypersetup{
  colorlinks=true,
  linkcolor=blue!55!black,
  urlcolor=blue!55!black,
  pdftitle={Free-field computation of the two-Virasoro branching coefficient},
  pdfauthor={Machine working note}
}

\newcommand{\Ff}{\mathsf F}
\newcommand{\SCA}{\mathsf{SCA}}
\newcommand{\Vir}{\mathsf{Vir}}
\newcommand{\V}{\mathcal V}
\newcommand{\M}{\mathcal M}
\newcommand{\one}{\mathbf 1}
\newcommand{\se}{s_{\mathrm{even}}}
\newcommand{\so}{s_{\mathrm{odd}}}
\newcommand{\hatrho}{\widehat\rho}

\title{Free-Field Computation of the
\(\Vir\oplus\Vir\) Branching Coefficient\\[0.3em]
\large Highest-vector prescription, norm, and three-point numerator}
\author{Machine working note}
\date{\today}

\begin{document}

\maketitle

\begin{abstract}
The human note reduces the all-NS superconformal block to two ordinary
Virasoro blocks and leaves one quantity undetermined:
the normalized three-point coefficient of the
\(\Vir\oplus\Vir\) highest-weight vectors \(v_n\).
Here we compute that quantity from the free-field prescription of
arXiv:1111.2803.  Writing \(k=2n\in\mathbb Z\), the vector is a finite
Fermi-sea product of the modes
\(\chi_r=f_r-i\vartheta_r\).  Its three-point numerator is the factorized
blow-up factor \(l\) of that paper.  Specializing the same formula to the
identity operator gives the Shapovalov norm.  Combining the two produces a
closed product formula for the squared branching coefficient that appears in
the human note.  No Whittaker vector is used.
\end{abstract}

\tableofcontents

\section{NS free-field realization}

We first fix the realization in which the branching vector will be computed.
Set
\begin{equation}
 Q=b+b^{-1},\qquad
 c=\frac32+3Q^2,\qquad
 h(P)=\frac12\left(\frac{Q^2}{4}-P^2\right).
 \label{eq:free-field-parameters}
\end{equation}
All modes in this section are NS modes:
\(m\in\mathbb Z\) and \(r\in\mathbb Z+\frac12\).

\subsection{Oscillators and Fock vacuum}

The auxiliary Majorana fermion has modes \(f_r\); this is the fermion denoted
by \(\psi_r\) in the human note.  The free-field realization of the SCA uses
a Heisenberg oscillator \(c_m^\epsilon\), an NS fermion
\(\vartheta_r^\epsilon\), and a bosonic zero mode \(\mathcal P\).
Here \(\epsilon\in\{+1,-1\}\) labels two \emph{alternative reflected
realizations} of the same SCA module; the \(+\) and \(-\) oscillators are not
two simultaneous free fields.  Their nonzero brackets are
\begin{align}
 [c_m^\epsilon,c_n^\epsilon]&=m\delta_{m+n,0},
 &\{\vartheta_r^\epsilon,\vartheta_s^\epsilon\}&=\delta_{r+s,0},
 &\{f_r,f_s\}&=\delta_{r+s,0},
 \label{eq:free-field-algebra}\\
 [\mathcal P,c_m^\epsilon]&=[\mathcal P,\vartheta_r^\epsilon]=0,
 &[f_r,c_m^\epsilon]&=0,
 &\{f_r,\vartheta_s^\epsilon\}&=0.
 \notag
\end{align}
The auxiliary BPZ convention inherited from the human note is
\begin{equation}
 f_r^\dagger=-f_{-r}.
 \label{eq:f-BPZ}
\end{equation}

Let \(\one_{\Ff}\) be annihilated by \(f_r\) for \(r>0\), and let
\(\lvert P\rangle_\epsilon\) be annihilated by \(c_m^\epsilon\) and
\(\vartheta_r^\epsilon\) for \(m,r>0\), with
\(\mathcal P\lvert P\rangle_\epsilon=P\lvert P\rangle_\epsilon\).  We write
\begin{equation}
 \lvert\mathrm{vac};P\rangle
 :=\one_{\Ff}\otimes\lvert P\rangle_\epsilon.
 \label{eq:free-field-vacuum}
\end{equation}
Thus the tensor-product module is the span of arbitrary finite products of
the negative modes \(f_{-r}\), \(c_{-m}^\epsilon\), and
\(\vartheta_{-r}^\epsilon\) acting on
\(\lvert\mathrm{vac};P\rangle\).

\subsection{SCA generators}

In the realization labelled by \(\epsilon\), the abstract SCA generators act
for \(m\ne0\) as
\begin{align}
 L_m
 &=\frac12\sum_{j\ne0,m}c_j^\epsilon c_{m-j}^\epsilon
  +\frac12\sum_{r\in\mathbb Z+\frac12}
       \left(r-\frac m2\right)
       \vartheta_{m-r}^\epsilon\vartheta_r^\epsilon
  +\frac i2(Qm+2\epsilon\mathcal P)c_m^\epsilon,
 \notag\\
 G_r
 &=\sum_{m\ne0}c_m^\epsilon\vartheta_{r-m}^\epsilon
   +i(Qr+\epsilon\mathcal P)\vartheta_r^\epsilon,
 \label{eq:free-field-SCA}\\
 L_0
 &=\sum_{m>0}c_{-m}^\epsilon c_m^\epsilon
   +\sum_{r>0}r\vartheta_{-r}^\epsilon\vartheta_r^\epsilon
   +\frac12\left(\frac{Q^2}{4}-\mathcal P^2\right).
 \notag
\end{align}
These formulas obey the \(\mathcal N=1\) SCA relations at the central charge
in Eq.~\eqref{eq:free-field-parameters}, and
\(L_0\lvert P\rangle_\epsilon=h(P)\lvert P\rangle_\epsilon\).
The realization \(\epsilon=-1\) is used for \(k>0\), and
\(\epsilon=+1\) for \(k<0\).  They are related by the super-Liouville
reflection operator.

It is now meaningful to define the isotropic fermionic modes
\begin{equation}
 \chi_r^\epsilon:=f_r-i\vartheta_r^\epsilon,
 \qquad
 \{\chi_r^\epsilon,\chi_s^\epsilon\}=0.
 \label{eq:chi-definition}
\end{equation}
The vanishing anticommutator is the reason that the highest vector below is a
finite Fermi-sea product.

\subsection{The commuting Virasoro pair}

For completeness, define the auxiliary-fermion stress tensor and the mixed
bilinear
\begin{equation}
 L_n^{\Ff}=\frac12\sum_{r\in\mathbb Z+\frac12}
 r:f_{n-r}f_r:,
 \qquad
 U_n=\sum_{r\in\mathbb Z+\frac12}f_{n-r}G_r.
 \label{eq:LF-U}
\end{equation}
The two commuting Virasoro generators used in the human note are
\begin{align}
 L_n^{(1)}
 &=\frac{b^{-1}}{b^{-1}-b}L_n
   -\frac{b^{-1}+2b}{b^{-1}-b}L_n^{\Ff}
   +\frac{1}{b^{-1}-b}U_n,
 \notag\\
 L_n^{(2)}
 &=\frac{b}{b-b^{-1}}L_n
   -\frac{b+2b^{-1}}{b-b^{-1}}L_n^{\Ff}
   +\frac{1}{b-b^{-1}}U_n.
 \label{eq:two-Vir-generators}
\end{align}
They obey
\begin{equation}
 [L_m^{(1)},L_n^{(2)}]=0,
 \qquad
 c^{(i)}=1+6\bigl(Q^{(i)}\bigr)^2,
 \qquad
 Q^{(i)}=b^{(i)}+\bigl(b^{(i)}\bigr)^{-1},
 \label{eq:two-Vir-central-charges}
\end{equation}
where
\begin{equation}
 b^{(1)}=\frac{2b}{\sqrt{2-2b^2}},
 \qquad
 \bigl(b^{(2)}\bigr)^{-1}
 =\frac{2b^{-1}}{\sqrt{2-2b^{-2}}}.
 \label{eq:two-Vir-couplings}
\end{equation}
This finishes the algebraic setup.  Every branching vector computed below is
an explicit vector in the Fock space of
Eqs.~\eqref{eq:free-field-algebra}--\eqref{eq:free-field-vacuum}.

\section{The quantity left open in the human note}

The human note uses
\begin{equation}
 \M_{\mathrm{NS}}(h)
 =\V_{\Ff_{1/2}}\otimes\V_{\SCA_{1/2}}(h),
 \qquad
 h=\frac12\left(\frac{Q^2}{4}-P^2\right),
 \qquad Q=b+b^{-1},
 \label{eq:module}
\end{equation}
with the auxiliary free-fermion factor first.  The branching rule is
\begin{equation}
 \M_{\mathrm{NS}}(h)
 \cong
 \bigoplus_{n\in\frac12\mathbb Z}
 \V_{c^{(1)}}(h_n^{(1)})\otimes
 \V_{c^{(2)}}(h_n^{(2)}),
 \qquad
 h_n^{(1)}+h_n^{(2)}=h+2n^2.
 \label{eq:branching}
\end{equation}
The corresponding highest-weight vector is denoted by \(v_n\), and its
parity is \(2n\pmod2\).

For an ordered trinion, the human note defines
\begin{equation}
 B_a(h_3,h_2,h_1;n_3,n_2,n_1)
 =\frac{
  \hatrho_a(v_{3,n_3},v_{2,n_2},v_{1,n_1})
 }{
  \lVert v_{3,n_3}\rVert
  \lVert v_{2,n_2}\rVert
  \lVert v_{1,n_1}\rVert
 },
 \qquad
 a=2(n_1+n_2+n_3)\pmod2.
 \label{eq:B-definition}
\end{equation}
Only \(B_a^2\) enters the conformal-block decomposition.  Our target is
therefore
\begin{equation}
 \mathcal B_a:=B_a^2
 =\frac{T_a^2}{N_{k_3}(P_3)N_{k_2}(P_2)N_{k_1}(P_1)},
 \qquad k_i:=2n_i,
 \label{eq:B-square-target}
\end{equation}
where
\begin{equation}
 T_a:=\hatrho_a(v_{3,n_3},v_{2,n_2},v_{1,n_1}),
 \qquad
 N_k(P):=\langle P,k\mid P,k\rangle.
 \label{eq:T-N}
\end{equation}

The order in Eq.~\eqref{eq:B-definition} is bra, inserted field, ket.  Other
slot orders are obtained with the Mobius and Koszul conventions already fixed
in the human note.

\section{Highest vector from the free-field realization}

\subsection{The Fermi-sea vector}

Set \(k=2n\in\mathbb Z\).  For \(k\ne0\), the prescription is
\begin{equation}
 |P,k\rangle
 =\Omega_k(P)
 \prod_{j=1}^{|k|}
 \chi_{-(2j-1)/2}^{-\operatorname{sgn}(k)}
 |\mathrm{vac};P\rangle,
 \qquad
 |P,0\rangle=|\mathrm{vac};P\rangle.
 \label{eq:free-field-vector}
\end{equation}
The product is ordered with increasing \(j\) from left to right.  Its total
level is
\begin{equation}
 \sum_{j=1}^{|k|}\frac{2j-1}{2}=\frac{k^2}{2}=2n^2,
 \label{eq:Fermi-level}
\end{equation}
which is precisely the shift in Eq.~\eqref{eq:branching}.

The factor \(\Omega_k(P)\) is fixed by the standard SCA-PBW normalization
\begin{equation}
 |P,k\rangle
 =\left(G_{-1/2}^{k^2}+\text{terms of lower degree in }G\right)|P\rangle.
 \label{eq:PBW-normalization}
\end{equation}
This fixes the vector, including its scale, without referring to a Whittaker
state.

The commutators quoted in the literature are
\begin{align}
 [L_m^{(1)}+L_m^{(2)},\chi_r^\epsilon]
 &=-\left(\frac m2+r\right)\chi_{m+r}^\epsilon,
 \notag\\
 [bL_m^{(1)}+b^{-1}L_m^{(2)},\chi_r^\epsilon]
 &=-\left((m+r)Q+\epsilon\mathcal P\right)
      \chi_{m+r}^\epsilon
   +i\sum_{j\ne0}c_j^\epsilon\chi_{m+r-j}^\epsilon.
 \label{eq:chi-commutators}
\end{align}
They prove directly that Eq.~\eqref{eq:free-field-vector} is killed by both
positive Virasoro algebras.  The two weights are
\begin{align}
 h_n^{(1)}
 &=\frac14(Q^{(1)})^2
  -\left(\frac{P}{\sqrt{2-2b^2}}+nb^{(1)}\right)^2,
 \notag\\
 h_n^{(2)}
 &=\frac14(Q^{(2)})^2
  -\left(\frac{P}{\sqrt{2-2b^{-2}}}+n(b^{(2)})^{-1}\right)^2.
 \label{eq:branch-weights}
\end{align}

The first nontrivial vectors, already converted back to the human note's SCA
and auxiliary-fermion modes, are
\begin{align}
 |P,1\rangle
 &=\left[G_{-1/2}+\left(\frac Q2+P\right)f_{-1/2}\right]|P,0\rangle,
 \notag\\
 |P,-1\rangle
 &=\left[G_{-1/2}+\left(\frac Q2-P\right)f_{-1/2}\right]|P,0\rangle.
 \label{eq:k-one-vectors}
\end{align}
They are the human-note vectors \(v_{1/2}\) and \(v_{-1/2}\).

\section{Factorized three-point numerator}

\subsection{Dictionary for an ordered trinion}

Parameterize the inserted NS weight by
\begin{equation}
 h_2=\frac12\alpha_2(Q-\alpha_2),
 \qquad
 \alpha_2=\frac Q2+P_2.
 \label{eq:alpha-dictionary}
\end{equation}
For the ordered matrix element
\begin{equation}
 \langle P_3,k_3|\mathbb V_{\alpha_2}^{(k_2)}(1)|P_1,k_1\rangle,
 \label{eq:ordered-matrix-element}
\end{equation}
the state-field correspondence identifies
\(\mathbb V_{\alpha_2}^{(k_2)}\) with the same highest vector
\(|P_2,k_2\rangle\) used on the middle leg.  The normalized ratio in
arXiv:1111.2803 is
\begin{equation}
 l(\alpha,m\mid P',k',P,k)
 =\begin{cases}
 \displaystyle
 \frac{\langle P',k'|\mathbb V_\alpha^{(m)}|P,k\rangle}
 {\langle P'|\mathbb V_\alpha^{(0)}|P\rangle},
 &m+k+k'\text{ even},\\[1.1em]
 \displaystyle
 \frac{\langle P',k'|\mathbb V_\alpha^{(m)}|P,k\rangle}
 {\langle P'|\mathbb V_\alpha^{(\pm1)}|P\rangle},
 &m+k+k'\text{ odd}.
 \end{cases}
 \label{eq:l-definition}
\end{equation}
The denominators are exactly the human note's kinematic normalizations:
\(\rho_0(\phi_3,\phi_2,\phi_1)=1\) in the even case and
\(\rho_1(\phi_3,G_{-1/2}\phi_2,\phi_1)=1\) in the odd case.  Therefore
\begin{equation}
 \boxed{
 T_a
 =l(\alpha_2,k_2\mid P_3,k_3,P_1,k_1),
 \qquad
 a=k_1+k_2+k_3\pmod2.}
 \label{eq:T-equals-l}
\end{equation}

\subsection{The product formula}

For \(\sigma,\tau\in\{+1,-1\}\), define
\begin{equation}
 x_{\sigma\tau}:=\alpha+\sigma P'+\tau P,
 \qquad
 r_{\sigma\tau}:=\frac{m+\sigma k'+\tau k}{2}.
 \label{eq:x-r}
\end{equation}
For a nonnegative integer \(r\), set
\begin{align}
 \se(x,r)
 &=2^{-r^2/2}
 \prod_{\substack{i,j\ge1,\ i+j\le2r\\i+j\equiv0\ ({\rm mod}\ 2)}}
 \left[x+(i-1)b+(j-1)b^{-1}\right],
 \notag\\
 \so(x,r)
 &=2^{-r(r+1)/2}
 \prod_{\substack{i,j\ge1,\ i+j\le2r+1\\i+j\equiv1\ ({\rm mod}\ 2)}}
 \left[x+(i-1)b+(j-1)b^{-1}\right].
 \label{eq:s-positive}
\end{align}
For negative integer \(r\), extend them by
\begin{equation}
 \se(x,r)=(-1)^r\se(Q-x,-r),
 \qquad
 \so(x,r)=\so(Q-x,-r).
 \label{eq:s-negative}
\end{equation}
Finally define
\begin{equation}
 \operatorname{Int}(x):=\operatorname{sgn}(x)\lfloor|x|\rfloor.
 \label{eq:Int}
\end{equation}

The numerator is then
\begin{equation}
 \boxed{
 l(\alpha,m\mid P',k',P,k)
 =\begin{cases}
 \displaystyle
 \prod_{\sigma,\tau=\pm1}
 \se\left(x_{\sigma\tau},r_{\sigma\tau}\right),
 &m+k+k'\text{ even},\\[1.2em]
 \displaystyle
 \prod_{\sigma,\tau=\pm1}
 \so\left(x_{\sigma\tau},
 \operatorname{Int}(r_{\sigma\tau})\right),
 &m+k+k'\text{ odd}.
 \end{cases}}
 \label{eq:l-product}
\end{equation}
This is already the desired unnormalized three-point coefficient of the
three highest vectors.

\section{Norm from the same formula}

Take the inserted field to be the identity:
\begin{equation}
 \alpha=0,
 \qquad m=0,
 \qquad P'=P,
 \qquad k'=k.
 \label{eq:norm-specialization}
\end{equation}
Then Eq.~\eqref{eq:l-definition} becomes the norm, while the four factors in
Eq.~\eqref{eq:l-product} reduce to
\begin{equation}
 (x_{++},r_{++})=(2P,k),
 \quad
 (x_{+-},r_{+-})=(0,0),
 \quad
 (x_{-+},r_{-+})=(0,0),
 \quad
 (x_{--},r_{--})=(-2P,-k).
 \label{eq:norm-four-factors}
\end{equation}
Since \(\se(x,0)=1\), we obtain
\begin{equation}
 \boxed{
 N_k(P)
 =\langle P,k|P,k\rangle
 =\se(2P,k)\se(-2P,-k).}
 \label{eq:norm-master}
\end{equation}
Equivalently, for \(p\ge0\),
\begin{align}
 N_p(P)&=(-1)^p\se(2P,p)\se(Q+2P,p),
 \notag\\
 N_{-p}(P)&=(-1)^p\se(Q-2P,p)\se(-2P,p).
 \label{eq:norm-piecewise}
\end{align}

At level \(1/2\), Eq.~\eqref{eq:k-one-vectors} gives an independent direct
check.  Using
\begin{equation}
 \langle G_{-1/2}\phi,G_{-1/2}\phi\rangle=2h,
 \qquad
 \langle f_{-1/2}\one,f_{-1/2}\one\rangle=-1,
 \label{eq:level-half-inner-products}
\end{equation}
and vanishing mixed inner product, one finds
\begin{align}
 N_1(P)
 &=2h-\left(\frac Q2+P\right)^2
 =-P(Q+2P),
 \notag\\
 N_{-1}(P)
 &=2h-\left(\frac Q2-P\right)^2
 =P(Q-2P).
 \label{eq:level-half-norms}
\end{align}
These agree with Eq.~\eqref{eq:norm-master}, because
\(\se(x,1)=x/\sqrt2\).  The same computation also gives
\begin{equation}
 \langle P,1|P,-1\rangle=0.
 \label{eq:level-half-orthogonality}
\end{equation}

\section{Closed formula for the human-note coefficient}

Combining Eqs.~\eqref{eq:B-square-target}, \eqref{eq:T-equals-l}, and
\eqref{eq:norm-master} gives
\begin{equation}
 \boxed{
 \mathcal B_a
 (P_3,P_2,P_1;k_3,k_2,k_1)
 =\frac{
 l\left(\frac Q2+P_2,k_2\mid P_3,k_3,P_1,k_1\right)^2
 }{
 N_{k_3}(P_3)N_{k_2}(P_2)N_{k_1}(P_1)
 },}
 \label{eq:B-final}
\end{equation}
where \(l\) is Eq.~\eqref{eq:l-product}, \(N\) is
Eq.~\eqref{eq:norm-master}, and
\begin{equation}
 k_i=2n_i,
 \qquad
 a=k_1+k_2+k_3\pmod2.
 \label{eq:final-dictionary}
\end{equation}
This is the coefficient that multiplies the product of the two Virasoro
blocks in the all-NS decomposition of the human note.

If an unsquared \(B_a\) is desired, one must choose square roots of the three
complex Shapovalov norms and fix an overall sign for every branch vector.
The sewn conformal block is insensitive to that choice and naturally uses
Eq.~\eqref{eq:B-final}.

\section{Low-level values}

The product formula immediately gives
\begin{equation}
 l(\alpha,0\mid P',0,P,0)=1,
 \qquad
 l(\alpha,1\mid P',0,P,0)=1.
 \label{eq:base-l}
\end{equation}
The first is the even primary normalization; the second is the odd
normalization.

With one \(k=1\) branch vector and the other two labels zero,
\begin{equation}
 l(\alpha,0\mid P',0,P,1)
 =l(\alpha,0\mid P',1,P,0)=1.
 \label{eq:one-branch-l}
\end{equation}
Consequently the corresponding squared normalized coefficient is simply the
inverse norm of that leg.

With two \(k=1\) vectors, the even numerators are
\begin{align}
 l(\alpha,0\mid P',1,P,1)
 &=\frac12(\alpha+P'+P)(\alpha-P'-P-Q),
 \notag\\
 l(\alpha,1\mid P',1,P,0)
 &=\frac12(\alpha+P'+P)(\alpha+P'-P),
 \notag\\
 l(\alpha,1\mid P',0,P,1)
 &=\frac12(\alpha+P'+P)(\alpha-P'+P).
 \label{eq:two-branch-l}
\end{align}
These are the first nontrivial polynomial checks to reproduce directly from
the human-note Ward identities.

For example,
\begin{equation}
 \mathcal B_0(P_3,P_2,P_1;1,0,1)
 =\frac{
  (\alpha_2+P_3+P_1)^2
  (\alpha_2-P_3-P_1-Q)^2
 }{
  4N_1(P_3)N_1(P_1)
 },
 \qquad \alpha_2=\frac Q2+P_2.
 \label{eq:B-example}
\end{equation}

\subsection{Complete unpacking for
\texorpdfstring{\(n_i=0,\pm\frac12\)}{ni = 0, +/- 1/2}}

We now remove the product notation completely at the first nontrivial
branching level.  In this subsection the slot order
\((n_1,n_2,n_3)\) is ket, inserted field, bra.  To keep the left-hand sides
readable, define only
\begin{equation}
 \mathcal B_a[n_1,n_2,n_3]
 :=\mathcal B_a(P_3,P_2,P_1;2n_3,2n_2,2n_1),
 \qquad
 T_a[n_1,n_2,n_3]
 :=\hatrho_a(v_{3,n_3},v_{2,n_2},v_{1,n_1}).
 \label{eq:small-n-slot-notation}
\end{equation}
No momentum argument is suppressed on the right-hand sides below.

Whenever \(n_i\ne0\), write
\begin{equation}
 \epsilon_i:=2n_i\in\{+1,-1\}.
 \label{eq:small-n-epsilon}
\end{equation}
The free-field vector and its norm take the uniform form
\begin{align}
 |P,\epsilon\rangle
 &=\left[
 G_{-1/2}+\left(\frac Q2+\epsilon P\right)f_{-1/2}
 \right]|P,0\rangle,
 \notag\\
 \langle P,\epsilon|P,\epsilon\rangle
 &=-\epsilon P\left(Q+2\epsilon P\right).
 \label{eq:small-n-vector-and-norm}
\end{align}

\paragraph{Zero or one nonzero label.}
For the vacuum branch,
\begin{equation}
 T_0[0,0,0]=1,
 \qquad
 \mathcal B_0[0,0,0]=1.
 \label{eq:small-n-vacuum}
\end{equation}
If exactly one label is nonzero, the numerator remains one and
\begin{align}
 T_1[\epsilon_1/2,0,0]
 &=T_1[0,\epsilon_2/2,0]
 =T_1[0,0,\epsilon_3/2]=1,
 \notag\\
 \mathcal B_1[\epsilon_1/2,0,0]
 &=-\frac{\epsilon_1}
 {P_1(Q+2\epsilon_1P_1)},
 \notag\\
 \mathcal B_1[0,\epsilon_2/2,0]
 &=-\frac{\epsilon_2}
 {P_2(Q+2\epsilon_2P_2)},
 \notag\\
 \mathcal B_1[0,0,\epsilon_3/2]
 &=-\frac{\epsilon_3}
 {P_3(Q+2\epsilon_3P_3)}.
 \label{eq:small-n-one-leg}
\end{align}
For example, a \(+\frac12\) label gives
\(-[P_i(Q+2P_i)]^{-1}\), whereas a \(-\frac12\) label gives
\([P_i(Q-2P_i)]^{-1}\).

\paragraph{Exactly two nonzero labels.}
When the ket label vanishes, one finds
\begin{align}
 T_0[0,\epsilon_2/2,\epsilon_3/2]
 &=\frac12\left[
 \left(\frac Q2+\epsilon_2P_2+\epsilon_3P_3\right)^2-P_1^2
 \right],
 \notag\\
 \mathcal B_0[0,\epsilon_2/2,\epsilon_3/2]
 &=
 \frac{
 \left[
 \left(\frac Q2+\epsilon_2P_2+\epsilon_3P_3\right)^2-P_1^2
 \right]^2
 }{
 \begin{gathered}
 4\epsilon_2\epsilon_3P_2P_3
 (Q+2\epsilon_2P_2)\\[-0.2em]
 {}\times(Q+2\epsilon_3P_3)
 \end{gathered}
 }.
 \label{eq:small-n-ket-zero}
\end{align}
When the inserted label vanishes, the ordered trinion gives instead
\begin{align}
 T_0[\epsilon_1/2,0,\epsilon_3/2]
 &=\frac12\left[
 P_2^2-
 \left(\frac Q2+\epsilon_1P_1+\epsilon_3P_3\right)^2
 \right],
 \notag\\
 \mathcal B_0[\epsilon_1/2,0,\epsilon_3/2]
 &=
 \frac{
 \left[
 P_2^2-
 \left(\frac Q2+\epsilon_1P_1+\epsilon_3P_3\right)^2
 \right]^2
 }{
 \begin{gathered}
 4\epsilon_1\epsilon_3P_1P_3
 (Q+2\epsilon_1P_1)\\[-0.2em]
 {}\times(Q+2\epsilon_3P_3)
 \end{gathered}
 }.
 \label{eq:small-n-insertion-zero}
\end{align}
Finally, when the bra label vanishes,
\begin{align}
 T_0[\epsilon_1/2,\epsilon_2/2,0]
 &=\frac12\left[
 \left(\frac Q2+\epsilon_1P_1+\epsilon_2P_2\right)^2-P_3^2
 \right],
 \notag\\
 \mathcal B_0[\epsilon_1/2,\epsilon_2/2,0]
 &=
 \frac{
 \left[
 \left(\frac Q2+\epsilon_1P_1+\epsilon_2P_2\right)^2-P_3^2
 \right]^2
 }{
 \begin{gathered}
 4\epsilon_1\epsilon_2P_1P_2
 (Q+2\epsilon_1P_1)\\[-0.2em]
 {}\times(Q+2\epsilon_2P_2)
 \end{gathered}
 }.
 \label{eq:small-n-bra-zero}
\end{align}
The relative sign in Eq.~\eqref{eq:small-n-insertion-zero} is part of the
ordered three-point prescription; it should not be replaced by a symmetric
guess.  It disappears after squaring the numerator but remains relevant if
an unsquared coefficient is assigned.

\paragraph{Three nonzero labels.}
All eight sign choices are covered by
\begin{align}
 T_1[\epsilon_1/2,\epsilon_2/2,\epsilon_3/2]
 &=\frac12
 \left(
 \frac Q2+\epsilon_1P_1+\epsilon_2P_2+\epsilon_3P_3+b
 \right)
 \notag\\[-0.2em]
 &\hspace{3.5em}\times
 \left(
 \frac Q2+\epsilon_1P_1+\epsilon_2P_2+\epsilon_3P_3+b^{-1}
 \right),
 \label{eq:small-n-three-leg-T}\\
 \mathcal B_1[\epsilon_1/2,\epsilon_2/2,\epsilon_3/2]
 &=
 -\frac{
 \begin{gathered}
 \epsilon_1\epsilon_2\epsilon_3
 \left(
 \frac Q2+\epsilon_1P_1+\epsilon_2P_2+\epsilon_3P_3+b
 \right)^2\\[-0.2em]
 {}\times
 \left(
 \frac Q2+\epsilon_1P_1+\epsilon_2P_2+\epsilon_3P_3+b^{-1}
 \right)^2
 \end{gathered}
 }{
 \begin{gathered}
 4P_1P_2P_3
 (Q+2\epsilon_1P_1)(Q+2\epsilon_2P_2)\\[-0.2em]
 {}\times(Q+2\epsilon_3P_3)
 \end{gathered}
 }.
 \label{eq:small-n-three-leg-B}
\end{align}
Equations~\eqref{eq:small-n-vacuum}--\eqref{eq:small-n-three-leg-B}
exhaust all \(27\) choices \(n_i\in\{0,\pm\frac12\}\).  They are precisely
the formulas evaluated by
\texttt{Code/two\_virasoro\_fusion.py}.

\section{What remains to be audited}

The calculation above settles the all-NS coefficient, modulo a direct audit
of the literature product formula against the human note's ordered Ward
identities.  The efficient audit is:
\begin{enumerate}
 \item reproduce Eqs.~\eqref{eq:level-half-norms} from the Gram form;
 \item reproduce Eqs.~\eqref{eq:one-branch-l} and
 Eq.~\eqref{eq:two-branch-l} from the three-point Ward identities;
 \item check a \(|k|=2\) case using the explicit level-two vectors in
 Appendix B of arXiv:1111.2803; and
 \item only then use Eq.~\eqref{eq:B-final} as the all-level input to the
 block decomposition.
\end{enumerate}

The Ramond branching vectors have an analogous free-field prescription in
arXiv:1608.02568.  It uses integer \(\chi_r^\pm\) modes and has two parity
copies.  Their norms and \((\mathrm{NS},\mathrm R,\mathrm R)\) three-point
coefficients require a separate translation to the human note's phase
convention
\begin{equation}
 G_0w_\beta^\pm=i\beta e^{\mp i\pi/4}w_\beta^\mp.
\end{equation}
No Whittaker coefficient is used here.

\section{Sources}

The free-field vector, its PBW normalization, the blow-up factors, and the
functions \(\se,\so\) are in
\href{https://arxiv.org/abs/1111.2803}{arXiv:1111.2803}, especially
Sec.~3.2, Proposition 3.2, and Appendix B.  The later presentation and the
Ramond analogue are in
\href{https://arxiv.org/abs/1608.02568}{arXiv:1608.02568}, Sec.~4.

\end{document}
